\subsection{Ergebnisse der Umsetzung}
Durch die Umsetzung war der Production Cluster eingerichtet und konnte Codetask ausführen. Diese Codetask Instance ist über VPN erreichbar.

Durch den Wechsel zur GitOps-Verwaltung ergeben sich einige Vorteile.
Einerseits gibt es eine erhöhte Sicherheit, da kein weiteres Tool wie z. B. die Pipeline Zugriff auf den Cluster benötigt, sondern der Cluster selbst durch FluxCD auf Änderungen reagiert.
Außerdem wurde die Verwendung des Clusters einfacher, da die Studierenden das Deployment mit Git-Befehlen verwalten können und nicht direkt Helm- und Kubernetes-Befehle erlernen müssen.
Die Cluster sind besser dokumentiert, da durch den deklarativen Ansatz die Infrastruktur und Anwendungen im Soll-Zustand stehen und Veränderungen durch die Commit-History geloggt werden.
Das Self-Healing hat sich durch die Drift-Detection verbessert, da regelmäßig geprüft wird, ob der Ist-Zustand im Cluster dem Soll-Zustand im Git-Repository entspricht. 
Das Deployment ist nun auch isoliert von der restlichen Anwendung und wird zentral verwaltet. Das Codetask-Repository ist nicht mehr für das Deployment verantwortlich, sondern muss in der Pipeline nur noch den Code testen und die Images bauen.

Aber auch durch die Erweiterungen von Longhorn und MetalLB können Single Points of Failure verhindert werden. Da MetalLB externe IP-Adressen verwendet, muss keine der Node-IP-Adressen exponiert werden.
Das Horizontal Pod Scaling verbessert die Stabilität, da die Last nun besser skaliert wird. 
Auch das gleichmäßige Verteilen der Pods auf die Nodes führt zu einer höheren Verfügbarkeit, Stabilität und Fehlertoleranz.